\documentclass[a4paper,14pt]{extarticle} % 14й шрифт
%%% Преамбула %%%

\usepackage{fontspec} % XeTeX
\usepackage{xunicode} % Unicode для XeTeX
\usepackage{xltxtra}  % Верхние и нижние индексы
\usepackage{pdfpages} % Вставка PDF
\usepackage{bm}

\usepackage{listings} % Оформление исходного кода
\lstset{
    basicstyle=\small\ttfamily, % Размер и тип шрифта
    breaklines=true, % Перенос строк
    tabsize=2, % Размер табуляции
    literate={--}{{-{}-}}2 % Корректно отображать двойной дефис
}
% Русский язык
\usepackage{polyglossia}
\setdefaultlanguage{russian}
\setotherlanguage{english}

% Шрифты, xelatex
\defaultfontfeatures{Ligatures=TeX}
\setmainfont[Ligatures=TeX]{Times New Roman}
\newfontfamily\cyrillicfont{Times New Roman}[Script=Cyrillic]
%\setsansfont{Liberation Sans} % Альтернатива нью роману, если его нет
\setmonofont{FreeMono} % Моноширинный шрифт для оформления кода
\newfontfamily{\cyrillicfonttt}{FreeMono}


\usepackage{amssymb,amsfonts,amsmath} % Математика
\numberwithin{equation}{section} % Формула вида секция.номер

\usepackage{enumerate} % Тонкая настройка списков
\usepackage{indentfirst} % Красная строка после заголовка
\usepackage{float} % Расширенное управление плавающими объектами
\usepackage{multirow} % Сложные таблицы

\usepackage{pdfcomment}     % тудущки и комменты
\pdfcommentsetup{final} % set to '' to prevent all comments from showing up   draft
\pdfcommentsetup{color={1.0 1.0 0.0}, open=true}

% Пути к каталогам с изображениями
\usepackage{graphicx} % Вставка картинок и дополнений
\graphicspath{{pics/}{images/userguide/}{images/testing/}{images/infrastructure/}{extra/}{extra/drafts/}}

% Формат подрисуночных записей
\usepackage{chngcntr}
\counterwithin{figure}{section}

% Гиперссылки
\usepackage{hyperref}
\hypersetup{
    colorlinks, urlcolor={black}, % Все ссылки черного цвета, кликабельные
    linkcolor={black}, citecolor={black}, filecolor={black},
    pdfauthor={Author},
    pdftitle={Title}
}
% Оформление библиографии и подрисуночных записей через точку
\makeatletter
\renewcommand*{\@biblabel}[1]{\hfill#1.}
\renewcommand*\l@section{\@dottedtocline{1}{1em}{1em}}
\renewcommand{\thefigure}{\thesection.\arabic{figure}} % Формат рисунка секция.номер
\renewcommand{\thetable}{\thesection.\arabic{table}} % Формат таблицы секция.номер
\def\redeflsection{\def\l@section{\@dottedtocline{1}{0em}{10em}}}
\makeatother

\renewcommand{\baselinestretch}{1.4} % Полуторный межстрочный интервал
\parindent 1.27cm % Абзацный отступ

\sloppy             % Избавляемся от переполнений
\hyphenpenalty=1000 % Частота переносов
\clubpenalty=10000  % Запрещаем разрыв страницы после первой строки абзаца
\widowpenalty=10000 % Запрещаем разрыв страницы после последней строки абзаца

% Отступы у страниц
\usepackage{geometry}
\geometry{left=3cm}
\geometry{right=1.5cm}
\geometry{top=2cm}
\geometry{bottom=2cm}

% Списки
\usepackage{enumitem}
\setlist[enumerate,itemize]{leftmargin=12.7mm} % Отступы в списках

\makeatletter
    \AddEnumerateCounter{\asbuk}{\@asbuk}{м)}
\makeatother
\setlist{nolistsep} % Нет отступов между пунктами списка
\renewcommand{\labelitemi}{--} % Маркет списка --
\renewcommand{\labelenumi}{\asbuk{enumi})} % Список второго уровня
\renewcommand{\labelenumii}{\arabic{enumii})} % Список третьего уровня
                    % Глубина оглавления, до subsubsection

% Нумерация страниц по центру снизу
\usepackage{fancyhdr}
\pagestyle{plain}
\fancyhf{}
\fancyhead[R]{\textrm{\thepage}}
\fancyheadoffset{0mm}
\fancyfootoffset{0mm}
\setlength{\headheight}{17pt}
\renewcommand{\headrulewidth}{0pt}
\renewcommand{\footrulewidth}{0pt}
\fancypagestyle{plain}{ 
    \fancyhf{}
    \rhead{\thepage}
}
% Содержание
\usepackage{tocloft}
\renewcommand{\cftaftertoctitle}{
    \thispagestyle{plain}}     % Убирает номер страницы
\renewcommand{\cfttoctitlefont}{\hspace{0.38\textwidth}\MakeTextUppercase} % СОДЕРЖАНИЕ
\renewcommand{\cftsecfont}{\hspace{0pt}}            % Имена секций в содержании не жирным шрифтом
\renewcommand\cftsecleader{\cftdotfill{\cftdotsep}} % Точки для секций в содержании
\renewcommand\cftsecpagefont{\mdseries}             % Номера страниц не жирные
\setcounter{tocdepth}{3}        

% GOST - biblatex
\usepackage[% 
bibstyle=gost-numeric, %подключение одного из четырех главных стилей biblatex-gost 
sorting=ntvy, %тип сортировки в библиографии
]{biblatex}
\addbibresource{refs.bib}

% Формат подрисуночных надписей
\RequirePackage{caption}
\DeclareCaptionLabelSeparator{defffis}{ -- } % Разделитель
\captionsetup[figure]{justification=centering, labelsep=defffis, format=plain, belowskip=-15pt} % Подпись рисунка по центру
\captionsetup[table]{justification=raggedright, labelsep=defffis, format=plain, singlelinecheck=false} % Подпись таблицы слева
\addto\captionsrussian{\renewcommand{\figurename}{Рисунок}} % Имя фигуры

% Пользовательские функции
\newcommand{\addimg}[4]{ % Добавление одного рисунка
    \begin{figure}
        \centering
        \includegraphics[width=#2\linewidth]{#1}
        \caption{#3} \label{#4}
    \end{figure}
}
\newcommand{\addimghere}[4]{ % Добавить рисунок непосредственно в это место
    \begin{figure}[H]
        \centering
        \includegraphics[width=#2\linewidth]{#1}
        \caption{#3} \label{#4}
    \end{figure}
}
\newcommand{\addtwoimghere}[6]{ % Вставка двух рисунков
    \begin{figure}[H]
        \centering
        \includegraphics[width=#2\linewidth]{#1}
        \hspace{1cm}
        \includegraphics[width=#4\linewidth]{#3}
        \caption{#5} \label{#6}
    \end{figure}
}
\newcommand{\addimgapp}[2]{ % Это костыль для приложения Б
    \begin{figure}[H]
        \centering
        \includegraphics[width=1\linewidth]{#1}
        \caption*{#2}
    \end{figure}
}

% Заголовки секций в оглавлении в верхнем регистре
\usepackage{textcase}
\makeatletter
\let\oldcontentsline\contentsline
\def\contentsline#1#2{
    \expandafter\ifx\csname l@#1\endcsname\l@section
        \expandafter\@firstoftwo
    \else
        \expandafter\@secondoftwo
    \fi
    {\oldcontentsline{#1}{#2}}          % {\oldcontentsline{#1}{\MakeTextUppercase{#2}}} 
    {\oldcontentsline{#1}{#2}}
}
\makeatother

% Оформление заголовков
\usepackage[compact,explicit]{titlesec}
\titleformat{\section}{\bfseries}{}{0mm}{\centering{\thesection\quad\textbf{#1}}}   % \titleformat{\section}{}{}{12.5mm}{\centering{\thesection\quad\MakeTextUppercase{#1}}\vspace{1.5em}}
\titleformat{\subsection}[block]{\bfseries}{}{12.5mm}{\thesubsection\quad\textbf{#1}}          %
\titleformat{\subsubsection}[block]{\vspace{1em}\normalsize}{}{12.5mm}{\thesubsubsection\quad#1\vspace{1em}}
\titleformat{\paragraph}[block]{\normalsize}{}{12.5mm}{\MakeTextUppercase{#1}}


% Секции без номеров (введение, заключение...), вместо section*{}
\newcommand{\anonsection}[1]{
    \phantomsection % Корректный переход по ссылкам в содержании
    \paragraph{\centerline{{#1}}\vspace{0em}}
    \addcontentsline{toc}{section}{\uppercase{#1}}
}

% Секции для приложений
\newcommand{\appsection}[1]{
    \phantomsection
    \paragraph{\centerline{{#1}}}
    \addcontentsline{toc}{section}{\uppercase{#1}}
}

% Библиография: отступы и межстрочный интервал
\makeatletter
\renewenvironment{thebibliography}[1]
    {\section*{\refname}
        \list{\@biblabel{\@arabic\c@enumiv}}
           {\settowidth\labelwidth{\@biblabel{#1}}
            \leftmargin\labelsep
            \itemindent 16.7mm
            \@openbib@code
            \usecounter{enumiv}
            \let\p@enumiv\@empty
            \renewcommand\theenumiv{\@arabic\c@enumiv}
        }
        \setlength{\itemsep}{0pt}
    }
\makeatother

\setcounter{page}{2} % Начало нумерации страниц % Подключаем преамбулу

%%% Начало документа
\begin{document}

\includepdf{include/titlepage}
%\includepdf{pz} % Пояснительная записка
%\includepdf[pages={1,2}]{task} % Задание на диплом 
\include{include/referat.tex}

\tableofcontents % Содержание 
\clearpage

\anonsection{введение}

За последние годы был совершён существенный прогресс в доступности и точности сенсоров, позволяющих мобильным роботам 
осуществлять оценку окружающего пространства. Такие информационно-измерительные устройства как лидары, сонары (и что-нибудь 
ещё) стали основной опорой алгоритмов для алгоритмов автономной навигации и локализации. Тем не менее в роботах по-прежнему 
присутствуют оптические системы, так как они дают наиболее читаемую информацию для оператора в случаях, когда его вмешательство 
необходимо. Существенная часть современных мобильных роботов имеют у себя на борту камеры с широким () или сверхшироким ()          %TODO: найти значения углов
углами обзора, так как они, хоть и вносят искажения в воспринимаемую картину, позволяют охватить больше окружающего пространства. %FIXME:  Формулировки. 
Набор таких камер может составлять систему кругового обзора \cite{}, позволяющую оператору видеть не только в любом направлении, 
но даже с видом от третьего лица \cite{}. Схема подобной системы приведена на рисунке \ref{pic:examples}. 

\begin{figure}[H]
    \begin{center}
        \includegraphics[scale=0.5]{pics/sample.png}
        \caption{Системы кругового обзора}
        \label{pic:examples}
    \end{center}
\end{figure}
    
В случае автономных мобильных роботов подобные системы включаются лишь по необходимости, но при этом могут быть весьма 
дорогостоящими и занимать место в корпусе. %TODO: уточнить про стоимость. Перефразировать предложение. ценное место в корпусе? 
Согласно схемам на рисунке \ref{pic:examples} широкоугольные камеры в системах кругового обзора роботов часто имеют области % FIXME: 'часто' - оценочный термин?
пересечения их полей зрения, что позволяет проводить оценку глубины / использовать алгоритмы стереозрения / использовать камеры как стереопару. % FIXME: выбрать хорошую формулировку
Это даёт роботу вспомогательный (или единственный) источник трёхмерной информации об окружении без дополнительных расходов. % FIXME: разве что на вычислитель, вероятно...

Однако значительные радиальные искажения изображения, вызванные особенностями используемых объективов, и тот факт, что области пересечения 
обычно расположены ближе к краям изображения, не позволяют использовать известные алгоритмы стереозрения. % FIXME:не совсем алгоритмы стереозрения, скорее калибровки. Дополнить абзац

Целью работы является разработка и изучение точности алгоритма стереозрения, работающего с изображениями, полученными с использованием
широкоугольных и сверхширокоугольных объективов. 

В ходе работы решаются следующие задачи:
\begin{itemize}
    \item Обзор современных алгоритмов стереозрения / алгоритмов калибровки изображений широкоугольных камеры
    \item ...
\end{itemize}

\section{Обзор опыта построения систем стереозрения с использованием сверхширокоугольных камер}
\subsection{Модель сверхширокоугольной камеры}
\label{camera_model}
Сверхширокоугольные объективы имеют в своей основе сложную систему линз, схема которой вместе с примером изображения представлена на рисунке \ref{pic:fyscheme}. 
Особенности этой системы позволяют достигать существенного угла обзора, но также являются причиной аберрации и характерных искажений 
изображения. Моделировать реальный ход лучей в подобных камерах нецелесообразно, поэтому исследователи прибегают к моделям камер.   % FIXME: meh

\addtwoimghere{fisheye_scheme}{0.5}{fisheye_example}{0.5}{Схема хода лучей объектива "рыбий глаз" (слева), пример изображения (справа)\ref{fy_exmp}}{pic:fyscheme}

Как видно по рисунку \ref{pic:fy_geom}, модель проекции для камеры это функция (обычно обозначаемая $\pi_c()$), которая моделирует преобразование 
из точки трёхмерного пространства ($P=[x_c, y_c, z_c]^T$) в области зрения камеры в точку на плоскости изображения ($p=[u, \nu]^T$). Единичная            % не совсем единичная 
полусфера $S$ с центром в точке $O_c$ описывает поле зрения. На ней также лежит точка $P_C$, являющаяся результатом обратной проекции $\pi^{-1}_c({p})$.
Угол $\theta$ является углом падения для рассматриваемой точки, а угол $\phi$ откладывается между положительным направлением оси $x$ и $O_{i}{p}$. 

\addimghere{projection_geometry}{0.5}{Схема проекции точки трёхмерного пространства в точку на изображении}{pic:fy_geom}

Модели камер включают в себя описания нескольких типов искажений, но в сверхширокоугольных объективах самыми существенными являются радиальные - искажения, 
проявляющиеся сильнее ближе к краям изображения. Поэтому далее в этой секции модели будут рассматриваться именно с точки зрения описания радиальных искажений. 

Перспективная проекция, которая обычно используется в качестве модели ортоскопической камеры, не способна спроецировать широкоугольное пространство на кадр 
конечного размера. Поэтому при описании и проектировании fisheye-объективов опираются на другие виды проекций \cite{projections}. Но реальные линзы не всегда в 
точности следуют заданным моделям, к тому же отличия в используемых параметрах усложняют процесс калибровки камер. По этой причине радиальные искажения принято 
аппроксимировать многочленами, например, вида
 \begin{equation}	\pdfcomment{ переписать уравнение }
	\begin{split}
        \delta r= k_1 r^3 + k_2 r^5 + k_3 r^7 + ... + k_n r^{n+2},
        \label{eqn:fisheye_distortion}
    \end{split}
\end{equation}
где $k_i$ - коэффициенты, описывающие внутренние параметры камеры. 

В настоящий момент есть несколько распространённых моделей, аппроксимирующих реальные искажения подобных объективов. Модель Канналы и 
Брандта \cite{opencv_model} реализована в OpenCV и описывает радиальные искажения через угол падения луча света на линзу, а не расстояние  
от центра изображения до места падения, как это делалось в более ранних моделях. Авторы посчитали, что для описания типичных искажений достаточно 
пяти членов полинома. Таким образом, указанную модель можно записать следующими уравнениями:
\begin{equation}	
        \theta = \arctan(\frac{r}{f}),
        \label{eqn:kannala_theta}
\end{equation}

\begin{equation}	
    \delta r = k_1\theta + k_2\theta^3 + k_3\theta^5 + k_4\theta^7 + k_5\theta^9,
    \label{eqn:kannala_r}
\end{equation}

\begin{equation}	
    \begin{pmatrix}u\\v\end{pmatrix} = \delta r(\theta)\begin{pmatrix}cos(\phi)\\sin(\phi)\end{pmatrix},
    %\delta r = k_1\theta + k_2\theta^3 + k_3\theta^5 + k_4\theta^7 + ... + k_n\theta^{n+1}
    \label{eqn:kannala_uv}
\end{equation}
где $\theta$ - угол падения луча, определяемый выбранным типом проекции; $\phi$ - угол между горизонтом 
и проекцией падающего луча на плоскость изображения; $r$ - расстояние от спроектированной точки до центра
изображения; $f$ - фокусное расстояние.

Также большое распространение получила модель Скарамуззы \cite{scaramuzza}, которая легла в основу Matlab Omnidirectional 
Camera Calibration Toolbox. Она связывает точки на изображении с соответствующей им точкой в координатах камеры 
следующим образом
\begin{equation}	
    \begin{pmatrix}X_c\\Y_c\\Z_c\end{pmatrix} = \lambda \begin{pmatrix}u\\v\\a_0 + a_2 r^2 + a_3 r^3 + a_4 r^4\end{pmatrix},
    %\delta r = k_1\theta + k_2\theta^3 + k_3\theta^5 + k_4\theta^7 + ... + k_n\theta^{n+1}
    \label{eqn:scaramuzza}
\end{equation}
где $a_0 ... a_4$ - коэффициенты, описывающие внутренние параметры камеры; $\lambda$ - масштабный коэффициент.



\pdfcomment{ дописать  например double sphere model }



\subsection{Обзор существующих систем стереозрения, использующих сверхширокоугольные камеры}

Распространённые библиотеки машинного зрения (OpenCV, MATLAB CV Toolbox) предлагаю готовые к использованию классы и функции, позволяющие после калибровки
камер получать с помощью них карты глубины. Однако на практике эти методы весьма ограничены. Для стереосопоставления 
используются традиционные методы, приспособленные для классических камер с перспективной проекцией, что не позволяет 
использовать кадры с широкоугольных камер целиком. В результате у этих кадров после устранения искажений остаётся 
угол зрения порядка $90^\circ$. Кроме того, библиотечные функции не позволяют задать область интереса для каждой камеры, 
что ограничивает их область применения только для копланарного расположения камер.   \pdfcomment{ найти источник }

Исследователи предложили несколько реализаций стереозрения, опирающихся на снимки со сверхширокоугольных 
объективов и лишённых недостатков распространённых решений.  
% Вообще, в этой работе и вся та же тема с поворотами есть 
Например, метод, предложенный в \cite{omni_stereo}, позволяет создать кольцевую область стереозрения с вертикальным  %Метод с разворотом камер на 180  ...
полем зрения $65^\circ$. Для этого используются две $245^\circ$ камеры, закреплённые на противоположных концах жёсткого стержня.  
Это позволяет достигнуть панорамного обзора глубины с качеством, достаточным для осуществления автономной навигации и
локализации \cite{omni_copter}, но доступна такая схема расположения камер только летательным аппаратам.  

Другие авторы \cite{direct_neuro_stereo} решили отказаться от типичных для стереозрения этапов устранения искажений и ректификации % FIXME: пояснить ректификацию? 
и извлекать информацию о глубине напрямую по двум снимкам fisheye-камер. Для производства карт глубины используется 
свёрточная неиронная сеть, что требует существенных вычислительных мощностей - для достижения производительности в реальном 
времени разработчикам понадобилось использовать компьютер с ЦПУ i7-4770 и ГПУ NVIDIA GTX 1080Ti. В мобильном автономном роботе
такой вычислитель разместить может быть проблематично. Кроме того, метод так же предполагает, что обе камеры направлены в одном 
направлении. 

\pdfcomment{ монокулярные fisheye-методы }

% Были предложены ортогональные методы, но до реализации не дошло ...
 
\subsection{Обоснование выбора ПО}
\label{sec:software}
Разработку и первоначальные испытания алгоритма стереозрения целесообразно проводить в виртуальной среде. Это позволяет значительно упростить разработку, 
так как уменьшает время на проверку гипотез и расходы на реальное оборудование, особенно в случае неудачных испытаний. Из-за этих факторов виртуальное моделирование 
в робототехнике приобрело широкое распространение и активно применяется, например, для разработки систем локализации и навигации беспилотного транспорта \cite{}. 
Возросшее качество компьютерной графики к тому же позволило моделировать реалистичное окружение, что особенно важно при работе с системами технического зрения. 

Для разработки алгоритма, описанного в этой работе, нужна виртуальная среда, в которой можно симулировать несколько широкоугольных камер и настраивать их параметры, 
легко интегрировать алгоритмы технического зрения и создать окружение, приближенное к тому, в котором будет работать алгоритм. На данный момент исследователю доступен 
широкий выбор программного обеспечения, подходящего для этой задачи. В таблице \ref{tab:sims} представлено сравнение имеющихся предложений по основным изложенным выше 
требованиям.

\begin{table}[H]             % TODO: Сыровато. Надо дополнить. Проверить ГОСТовость подписи  FIXME: Таблица не умещается, но вращать не хочется. 
    \caption{Сравнение ПО для симуляции }
    \label{tab:sims}
    \begin{tabular}{|l|l|l|l|l|l|}
    \hline
    \textbf{Название} & \textbf{\begin{tabular}[c]{@{}l@{}}Симуляция \\ fisheye-камер\end{tabular}} & \textbf{\begin{tabular}[c]{@{}l@{}}Реалистичное \\ моделирование\end{tabular}} & \textbf{Интеграция кода}                                            & \textbf{Доступность} & \textbf{Примечания} \\ \hline
    Gazebo            & Возможна                                                                    & Затруднено                                                                     & \begin{tabular}[c]{@{}l@{}}Возможна \\ посредством ROS\end{tabular} & Бесплатно            &                     \\ \hline
    RoboDK            & Нет                                                                         & Затруднено                                                                     & Нет                                                                 & От 145€              &                     \\ \hline
    Webots            & Затруднена                                                                  & Возможно                                                                       & Возможна                                                            & Бесплатно            &                     \\ \hline
    CoppeliaSim       & Затруднена                                                                  & Затруднено                                                                     & Возможна                                                            & Бесплатно            &                     \\ \hline
    NVIDIA Isaac Sim  & Возможна                                                                    & Возможно                                                                       & Возможна                                                            & Бесплатно            &                     \\ \hline
    CARLA             & Затруднена                                                                  & Возможно                                                                       & Возможна                                                            & Бесплатно            &                     \\ \hline
    Unity             & Возможна                                                                    & Возможно                                                                       & Возможна                                                            & Бесплатно            &                     \\ \hline
                      &                                                                             &                                                                                &                                                                     &                      &                     \\ \hline
    \end{tabular}
\end{table}

По результатам оценки собранные сведений принято решение проводить разработку в симуляторе Unity. Он позволяет подробно настраивать камеру и эмулировать fisheye-объектив, 
строить реалистичные сцены благодаря свободному импорту моделей, а при программировании в симуляторе можно использовать сторонние программы в виде динамически подключаемых библиотек. 
По функционалу так же подходит NVIDIA Isaac Sim, но от него пришлось отказаться из-за высоких системных требований и малой изученности продукта.     % TODO: звучит тупо

Разрабатываемое решение должно иметь возможность внедрения в ПО робота, % TODO: не совсем ПО всё-таки
поэтому должно реализовываться на одном из популярных и быстродейственных языков программирования. Учитывая необходимость интеграции с Unity и потребность использовать популярные 
библиотеки, выбран язык C++.  Другим важным фактором является библиотека обработки изображений. В качестве основы для программной части была выбрана библиотека OpenCV, являющаяся стандартом 
при разработке систем технического зрения. Она доступна к использованию со множеством языков программирования, но наилучшую производительность показывает именно с C++ \cite{}.                              % мб есть сравнение
\subsection{Выводы по главе}

\newpage
\section{Система стереозрения}
\subsection{Алгоритм устранения искажений}

Как уже упоминалось в секции \ref{camera_model}, существующие способы устранения радиальных искажений не
позволяют работать близко к краям изображений, поэтому для реализации предлагаемой системы стереозрения     % FIXME: "работать"
был разработан алгоритм устранения искажений на основе модели \cite{scaramuzza}. Схема геометрического
принципа, лежащего в основе этого алгоритма, представлена на рисунке \ref{pic:sweeping}.

\addimghere{projection_sweeping}{0.7}{Схема принципа устранения искажений}{pic:sweeping} 

Цель алгоритма - найти, куда на выходном изображении проектируются все пиксели из выбранного участка входного 
сверхширокоугольного изображения. Сделать это можно, выполнив обратное преобразование \ref{eqn:scaramuzza} для каждого пикселя 
fisheye-снимка и затем прямое преобразование модели камеры-обскуры для получения результирующей проекции. Однако такой подход 
приведёт к возникновению дефектов из-за несовпадения частот дискретизации двух изображений.                 % TODO:  скорее всего наврал. Надо проконсультироваться 
Поэтому вместо этого алгоритм выполняет обратное преобразование  для каждого пикселя итогового изображения $\nu$, 
находя таким образом соответствующую ему точку в системе координат камеры $({X_c, Y_c, Z_c})$ 
\begin{equation}
    \label{eq:uv_to_xyz}
    \left[\begin{matrix}x_c\\y_c\\z_c\\\end{matrix}\right] = \left[\begin{matrix} {(u+c_x)*z_c}/f \\  {(v+c_y)*z_c}/f \\ z_c \\\end{matrix}\right],
\end{equation} 
где $c_x, c_y$ - координаты центра изображения; $f$ - фокусное расстояние. 

Набор таких точек формирует прямоугольную область $\nu_p$ с центром в точке $O_p$ и является виртуальной камерой-обскурой        % FIXME: прямоугольник?
с оптической осью $Z_p$. Поворот точек, входящих в $\nu_p$, с помощью матрицы вращения $\bm{R}$ образует $\nu'_p$ и
 позволяет таким образом задать направление обзора и ориентацию виртуальной камеры. 
\begin{equation}
    \label{eq:sweeped}
    \left[\begin{matrix}x'_p\\y'_p\\z'_p\\\end{matrix}\right] = \left[\begin{matrix}x_p\\y_p\\z_p\\\end{matrix}\right] \bm{R}.
\end{equation}  
\begin{equation}
    \label{eq:R}
    \bm{R} = \left[\begin{matrix}\cos{\alpha}&-\sin{\alpha}&0\\\sin{\alpha}&\cos{\alpha}&0\\0&0&1\\\end{matrix}\right]\left[\begin{matrix}1&0&0\\0&\cos{\beta}&\sin{\beta}\\0&-\sin{\beta}&\cos{\beta}\\\end{matrix}\right]\left[\begin{matrix}\cos{\gamma}&0&-\sin{\gamma}\\0&1&0\\\sin{\gamma}&0&\cos{\gamma}\\\end{matrix}\right],
\end{equation} 
где $\alpha, \beta, \gamma$ - углы Эйлера. % TODO: точно ли Эйлера? 

Тогда обратная fisheye-проекция точек из $\nu'_p$ позволяют получить область $\nu_i$ исходного изображения с искомыми пикселями. 
Таким образом, зная как геометрически проектируется каждая точка из $\nu$ в $\nu_i$, можно перенести информацию о цвете и получить 
изображение с устранёнными радиальными искажениями в любой части поля зрения камеры. 

Так как в используемой модели искажения в точке считаются центрально симметричными и зависят только от её удаления от центра изображения, 
рассмотрим ход падающего луча в координатах $(Z_c, \rho)$, изображённый на рисунке \ref{pic:scara_graph}.

\addimghere{scara_graph}{0.5}{}{pic:scara_graph} 

Для модели (\ref{eqn:scaramuzza}) обратная проекция записывается как 
\begin{equation}
    \label{eq:back_scara}
    \left[\begin{matrix}u_i\\v_i\\\end{matrix}\right] = \left[\begin{matrix} \frac{x_c}{\lambda}  \\  \frac{y_c}{\lambda} \\\end{matrix}\right],
\end{equation}  
где $\lambda = \rho_c / \rho_i$ - масштабный коэффициент. 

Для нахождения $\rho_i$ был применён метод последовательных приближений. Блок-схема алгоритма приведена на рисунке \ref{pic:newton_scheme}.

\addimghere{sample}{0.5}{Блок-схема алгоритма поиска $\rho$}{pic:newton_scheme} 

Очевидно, весь процесс преобразования $\nu \rightarrow \nu_i$ требует выполнения существенного количества математических операций, что 
негативно сказывается на скорости, с которой алгоритм может обрабатывать изображения в реальном времени. Однако при неизменных параметрах 
модели алгоритм достаточно выполнить лишь один раз, записав результат в таблицу поиска - структуру данных, которая позволяет дальнейшие 
преобразования проводить по уже известным соотношениям между пикселями. Время построения таблицы поиска для выходного изображения 
разрешением 540*540 пикселей и входного изображения 1080*1080 пикселей составляет $\sim 2.2c$. % TODO: нормально замерить. Возможно посчитать O(.)  O(n), n - кол-во пикселей? 
Это время можно уменьшить, применив различные оптимизации и параллельные вычисления (проекция каждого пикселя может рассчитывать независимо). % FIXME: оптимизации это про симметричность преобразований например. Надо нормально написать
Время отображения с использованием таблиц поиска составляет для того же изображения $ ... $, что позволяет применять алгоритм для устранения 
искажений в реальном времени.        % Как происходит выпрямление
\subsection{Описание системы стереозрения}

Предлагаемая система стереозрения состоит из двух камер с объективами типа "рыбий глаз" $\geqslant 180^\circ$,
расположенных ортогонально. В таком случае поля зрения соседних камер могут пересекаться, 
образуя  области, точки пространства в которых видны на обеих камерах. Это позволяет получить информацию о 
глубине в указанных областях. % FIXME: Всё описание просто бггг

Робот-доставщик "Ровер R3"  компании Яндекс имеет на борту 4 сверхширокоугольные камеры, размещённых спереди, сзади и по бортам 
корпуса \cite{yandex_rover}, что позволяет рассмотреть принцип работы системы стереозрения на его примере. На рисунке \ref{pic:4cam_system} 
представлено сравнение систем стереозрения. Слева изображена схема зон (обозначены зелёным), в которых обеспечивается получение информации о глубине
 при использовании описываемой системы, а справа эквивалентная по горизонтальному  покрытию схема при использовании обычных камер 
 с углом обзора $90^\circ$. 
 
%\addtwoimghere{Group 1}{0.4}{Group 2}{0.4}{Сравнение систем стереозрения}{pic:4cam_system}  % TODO: разобраться с масштабом
\addimghere{group12}{0.7}{Геометрическая модель бинокулярной системы стереозрения}{pic:2cam_scheme}

Как можно заметить, системе на основе обычных камер нужно в два раза больше сенсоров, чтобы достичь той же зоны покрытия 
по горизонтали.  Кроме того, традиционная система имеет меньшую зону покрытия по вертикали и не обеспечивает полный 
панорамный обзор. Все эти факторы делают систему стереозрения на основе ортогонально расположенных сверхширокоугольных камер 
более предпочтительной для применения в робототехнике.  % FIXME: уточнить. не во всей не всегда 

Применение традиционных алгоритмов стереозрения предполагает наличие стереопары - двух камер с известным взаимным положением,   % FIXME: meh как-то абзац теперь не на своём месте будто. Переформулировать наверное стоит
наблюдающих  одну область пространства с разных ракурсов. Описанный в предыдущей секции метод  устранения искажений 
сверхширокоугольной камеры позволяет создавать виртуальные камеры-обскуры и регулировать их направление обзора. Таким образом,
стереопару можно сформировать из двух таких виртуальных камер.  %параллельно биссектрисе угла, образованного пересечением полей зрения fisheye-камер. 
Далее для упрощения рассмотрения системы будет считаться, что оптические оси всех камер находятся в одной плоскости, 
а на ориентацию виртуальных камер влияет только один угол. % FIXME: Какой? ...только угол в этой плоскости? ...угол рыскания? 

На рисунке \ref{pic:2cam_scheme} изображён вариант системы с двумя камерами под углом $90^\circ$, соответствующий фрагменту схемы,     
  представленной на рисунке \ref{pic:4cam_system}. 

\addimghere{sample_simple2cam}{0.7}{Геометрическая модель бинокулярной системы стереозрения}{pic:2cam_scheme} %TODO: перерисовать схему?

  Здесь у полей зрения камер $C_0$ и $C_1$ есть область пересечения, обозначенная красным. 
Эта область эквивалентна области пересечения полей зрения двух камер с полями зрения $90^\circ$ (обозначены
оранжевым), повёрнутых на $45^\circ$ в сторону области интереса. В таком случае $B$ - база стереопары.
Примеры изображений, полученных в такой конфигурации, изображены на рисунке \ref{pic:dewarped_exmples}. 

\addimghere{4pic_example}{0.7}{Пример исходных изображений с отмеченной областью интереса и снимков виртуальной стереопары}{pic:dewarped_exmples}

      % Описание системы: Какие этапы, элементы и тд
\subsection{Виртуальное моделирование системы}
Описанная система была смоделирована в среде Unity, её внешний вид представлен на рисунке \ref{pic:unity_model}. Мир 
Unity предназначен для базовой проверки работоспособности испытываемого принципа, поэтому не содержит подробной модели
какого-либо робота. В нём присутствуют: плоскость земли, кронштейн, на котором сверхширокоугольные камеры закреплены под 
углом $90^\circ$, подвижный калибровочный узор и объекты-цели, предназначенные для оценки расстояния. В силу особенностей 
работы многих алгоритмов стереозрения эти объекты должны сильнотекстурированы \cite{disparity_review}. 

\addimghere{unity_view}{0.7}{Внешний вид сцены в Unity}{pic:unity_model} 

В Unity создание сцены происходит с использованием встроенных примитивов, импортированных файлов моделей популярных форматов
или моделей из магазина ассетов, предлагающего обширную библиотеку объектов и текстур, повторяющих различные реальные объекты. В данной 
работе использовались как и примитивы для создания простых объектов, так и модели из магазина для имитации препятствий и окружения.

Для моделирования камеры "рыбий глаз" использовался аддон Dome Tools из магазина Unity Asset Store.    
Он позволяет моделировать сверхширокоугольные объективы с разным углом зрения в эквидистантной проекции \cite{dome_tools}. 
При этом с точки зрения других искажений, не относящихся к моделированию правильной проекции, снимки с этой камеры получаются идеальными. % FIXME: моделированию чего? 
Окно настроек камеры представлено на рисунке \ref{pic:camera_settings}.

\addimghere{camera_settings}{0.5}{Окно настроек виртуальной камеры в Unity}{pic:camera_settings} 

Для виртуальной камеры доступны настройки угла зрения и виньетки по краям изображения,  % FIXME: виньетки -> абберации ??
а также различных параметров рендеринга, влияющих на уровень  детализации получаемого изображения.       % Моделирование системы 
\subsection{Передача изображений с виртуальных камер для обработки}

Unity в качестве основного языка программирования использует C\# и не позволяет импортировать 
выбранную библиотеку технического зрения OpenCV напрямую. Однако как уже упоминалось в секции 
\ref{sec:software}, данная среда позволяет интегрировать плагины, в том числе написанные на 
других языках программирования, в виде динамически подключаемых библиотек (DLL). Именно этот 
подход был использован для передачи изображений - была разработана библиотека с функциями для 
обработки снимков, которая компилировалась отдельно и затем импортировалась в Unity. 

Для реализации описанного в этой работе принципа и упрощения разработки в библиотеке были реализованы 
следующие функции:
\begin{itemize}
    \item initialize - выделяет память под нужные структуры, заполняет таблицу поиска и создаёт окна 
    для отображения будущих изображений.
    \item getImages - передаёт изображения в библиотеку и производит обработку изображений в 
    соответствии с аргументами. 
    \item takeScreenshot - производит ту же обработку входного изображения, что и предыдущая функция,
     но результат сохраняет в файл. 
    \item processImage - передаёт обратно в скрипт уже обработанные изображения для удобного отображения
    в интерфейсе Unity. 
    \item terminate - высвобождает память по завершении работы симуляции.      
\end{itemize}

Обработка, упомянутая в пунктах с функциями getImages и takeScreenshot, заключается в конвертации 
цветового пространства изображений (из RGBA, используемого в Unity, в BGR, используемый в OpenCV), 
горизонтальном зеркальном отображении (из-за разного положения точки отсчёта координат пикселей)
и, наконец, устранениb искажений в области интереса. Так как для наилучшей работы стереосопоставления
снимки должны быть синхронизированы (то есть получены камерами в один момент времени), обработка        % FIXME: ну не совсем одновременно и мб лучше это описать с другого ключа - типа функция может обрабатывать n снимков?
производится сразу для двух изображений. 

Также написан специальный скрипт, который управляет процессом взаимодействия симуляции и библиотеки. 
Эта программа имеет окно настроек, которое позволяет задать
используемые камеры, элементы управления параметрами  и выбрать область интереса с помощью задания углов.
 Внешний вид этого окна представлен на рисунке \ref{pic:connec_inter}.

\addimghere{connector_settings}{0.7}{Внешний вид окна настроек скрипта}{pic:connec_inter}

Сразу после запуска симуляции происходит запуск дополнительных вычислительных потоков для обработки 
кадров с каждой пары камер и передача настроек скриптом в библиотеку для построения таблиц поиска.           % FIXME: ну написано очень криво. Использование
Использование нескольких потоков позволяет не блокировать работу симуляции на время обработки изображений.  
Далее каждое обновление кадра скрипт считывает изображения с камер и помещает их в память соответствующего 
потока. Поток же работает независимо и обрабатывает каждую следующую пару изображений после готовности 
предыдущей.              % Про виртуальную камеру 
\subsection{Выводы по главе}       % Выводы по главе

\newpage
\section{Экспериментальное исследование системы стереозрения}
Описанная в предыдущей главе виртуальная модель системы стереозрения позволяет проводить 
с ней испытания для оценки работоспособности и сравнения с аналогами в контролируемой среде.            % FIXME: как же плохо       % План экспериментов

\subsection{Оценка качества устранения искажений}

Добавление в виртуальную модель эталонной камеры позволяет оценить качество устранения искажений путём сравнения
её снимков со снимками виртуальной камеры-обскуры модуля устранения искажений. % FIXME: никакой "модуль..." ранее не вводился, нужно переписать

Параметры полинома (\ref{eqn:scaramuzza}) для устранения искажений получены с помощью MATLAB Camera Calibrator и занесены 
в код библиотеки обработки изображений. % FIXME: опять-таки будто про OpenCV речь, надо как-то ссылаться на свой алгоритм\модуль
Снимки после устранения искажений и с эталонной камеры показаны на рисунке \ref{pic:2images_compar}.

\addtwoimghere{fy_house}{0.4}{reg_house}{0.4}{Слева - снимок после устранения искажений; справа - эталонный снимок.}{pic:2images_compar}

На левом изображении заметна большая резкость, вызванная, вероятно, округлениями чисел в процессе проекции. Также присутствует 
затенение по левому краю - следствие аберрации на исходном снимке. В той же области есть малозаметные искажения геометрии. Более 
явно эти дефекты можно увидеть на разностном изображении, представленном на рисунке \ref{pic:difference}. 

\addimghere{difference}{0.5}{Разностное изображение}{pic:difference}

Анализ этого изображения подтверждает различия в резкости и наличие искажений в левой части снимка. Тем не менее, искажения
достаточно несущественны и проявляются лишь близь краёв исходного изображения, что позволяет считать подобные снимки пригодными 
к применению в системе стереозрения.  Кроме того, некоторые эффекты возможно устранить или  сильно ослабить более качественной 
калибровкой камер... % TODO: и...  использованием более качественной оптики? применением системы не так близко к краям? 
         % Сравнение картинок

\subsection{С чем сравниваем}  % понятие облака точек не введено, поверхность какая? 

Результатом работы алгоритма стереозрения является карта глубины воспринимаемого пространства. Подобные карты могут
использоваться, например, алгоритмами навигации и локализации для построения карты окружения робота. Точность работы 
этих алгоритмов зависит от того, насколько точно карта глубины передаёт реальную информацию о форме объектов. Оценить 
эту характеристику для отдельной системы стереозрения проблематично, так как она зависит от множества факторов...
Поэтому оценка пригодности для применения системы в реальности проведена в сравнении с виртуальной моделью традиционной
стереопары.

Виртуальная модель позволяет разместить сразу несколько объектов в одной точке пространства, таким образом возможно
в существующую модель в Unity добавить ещё 2 камеры, совпадающие по параметрам и положению с виртуальными камерами на 
рисунке \ref{pic:2cam_scheme}. Разрешение этих камер выбрано исходя из размеров проекции $\nu_i$ области интереса на широкоугольном
снимке. Для камер "рыбий глаз" с разрешением $1080*1080$ пикселей она составляет 287482 пикселей, что аналогично камере 
с разрешением $540*540$.  В результате получена эталонная стереопара для сравнения. % TODO: чувствуется какая-то недосказанность

\subsection{Оценка качества устранения искажений}

Не уверен, что это стоит выносить в отдельный раздел. Вписывается не очень, надо переписывать подводку
\\
\\
Добавление в виртуальную модель эталонной камеры позволяет оценить качество устранения искажений путём сравнения
её снимков со снимками виртуальной камеры-обскуры модуля устранения искажений. % FIXME: никакой "модуль..." ранее не вводился, нужно переписать

Параметры полинома (\ref{eqn:scaramuzza}) для устранения искажений получены с помощью MATLAB Camera Calibrator и занесены 
в код библиотеки обработки изображений. % FIXME: опять-таки будто про OpenCV речь, надо как-то ссылаться на свой алгоритм\модуль
Снимки после устранения искажений и с эталонной камеры показаны на рисунке \ref{pic:2images_compar}.

\addtwoimghere{fy_house}{0.4}{reg_house}{0.4}{Слева - снимок после устранения искажений; справа - эталонный снимок.}{pic:2images_compar}

На левом изображении заметна большая резкость, вызванная, вероятно, округлениями чисел в процессе проекции. Также присутствует 
затенение по левому краю - следствие аберрации на исходном снимке. В той же области есть малозаметные искажения геометрии. Более 
явно эти дефекты можно увидеть на разностном изображении, представленном на рисунке \ref{pic:difference}. 

\addimghere{difference}{0.5}{Разностное изображение}{pic:difference}

Анализ этого изображения подтверждает различия в резкости и наличие искажений в левой части снимка. Тем не менее, искажения
достаточно несущественны и проявляются лишь близь краёв исходного изображения, что позволяет считать подобные снимки пригодными 
к применению в системе стереозрения.  Кроме того, некоторые эффекты возможно устранить или  сильно ослабить более качественной 
калибровкой камер... % TODO: и...  использованием более качественной оптики? применением системы не так близко к краям? 

\subsection{Оценка отклонения облака точек от поверхности}  % понятие облака точек не введено, поверхность какая? 

В качестве целевой поверхности выбрана виртуальная плоскость с нанесённой на неё текстурой высоко разрешения. Размеры и
положение плоскости относительно камер известно с высокой точностью, что позволяет сравнить результаты стереореконструкции
с реальным положением целевого объекта и оценить ошибку. В качестве метрики оценки выбрано среднее квадратичное отклонение
положения точек от модели плоскости. 

Исследования проведены в MATLAB, снимки получены с помощью виртуальной модели. Сначала для обеих стереопар проводится
калибровка по снимкам узора шахматной доски с помощью методики \cite{stereo_calib}, что позволяет получить 
внутренние и внешние параметры камер. % FIXME: ну вообще они у нас как бы есть. Мб сравнить ещё оценки с реальными 
Далее происходит построение карты расхождений \cite{SGBM} и 3D-реконструкция сцены. Результирующее облако точек 
представлено на рисунке
\ref{pic:raw_pointcloud}.

\addimghere{pointcloud}{0.7}{Неочищенное облако точек}{pic:raw_pointcloud}
% TODO: сюда можно и картинку после обработки поинтклауда \addtwoimghere{}{}{}{}{}
По заданным параметрам строится модель целевой плоскости, а точки за пределами её окрестности отбрасываются. % FIXME: размеры окрестности
Затем скрипт перебирает все оставшиеся точки и рассчитывает среднее квадратичное отклонение. Код скрипта приведён
 в приложении \ref{app:matlab}.

Описанный алгоритм применён к снимкам целевой плоскости на разных расстояниях с эталонной и исследуемой стереопар.
Результаты приведены на графике \ref{plot:MSE_compar}.

\addimghere{error_plot}{1.0}{График среднеквадратичной ошибки в зависимости от расстояния для эталонной и исследуемой стереопар}{plot:MSE_compar}

Как видно из графика, ошибка оценки поверхности в исследуемой системе в среднем в два раза больше этого показателя для
традиционной стереопары. Меньшая точность исследуемой системы связана с влиянием факторов, описанных в предыдущей секци,
на алгоритм стереосопоставления.  Тем не менее она остаётся в пределах $3\%$ от реального расстояния.  % FIXME: хммммм, как-то уж больно хорошо. Или  плохо... 10 см на 5 метрах



         % Оценка равномерности поинтклауда стены
\subsection{Выводы по главе}

Исследованы снимки, полученные с помощью виртуальных камер. Выполнено их сравнение с эталонными изображениями,  
которое продемонстрировало пригодность снимков для использования в системах стереозрения.

В виртуальную среду добавлена традиционная стереопара для сравнения с разработанной системой. Проведён эксперимент
по оценке качества облака точек, полученного с помощью предлагаемого решения. Результаты можно считать удовлетворительными 
для систем подобного класса. % TODO: \: meh  какого класса? Почему можно считать? 


  % Выводы по главе

\newpage
\anonsection{Заключение}

В ходе работы рассмотрены существующие системы стереозрения, в особенности использующие информацию
со сверхширокоугольных камер. Изучены модели, описывающие подобные камеры. Рассмотрено ПО для виртуального
 моделирования робототехнических комплексов и систем технического зрения. Выбрано ПО для моделировани и 
разработки системы стереозрения, использующей объективы "рыбий глаз". 

Разработан и описан алгоритм устранения искажений в области интереса и его математическая модель.   % пафосно, конечно.  Так ли хорошо она описана?
Описана возможная конфигурация системы стереозрения, рассмотрен её принцип работы. Система 
 смоделирована в среде Unity. В процессе размещены и настроены камеры, подготовлены объекты для калибровки и 
дальнейших испытаний системы. Реализована передача изображений с виртуальных камер для обработки алгоритмами
компьютерного зрения. С помощью снимков из виртуальной модели исследовано качество устранения искажений и выполнено
 сравнение точности оценки глубины предлагаемой системой и варианта с обычными камерами. Результаты позволяют 
 перейти к исследованию работы системы с применением реальных камер.  % FIXME: "указывают на доработки" такое себе 

Дальнейшая работа будет сконцентрирована на оптимизации и повышении точности алгоритма устранения искажений,
разработке способа автоматической установки параметров системы при разных конфигурациях расположения камер. 
Эти этапы позволят исследовать применимость системы стереозрения при  разных схемах размещения камер и разной 
длине базы, а также рассмотреть качество построения карты и локализации с применением описанной 
системы сначала с виртуальным, а затем и с реальном роботом. %Заключение-заключение

\newpage
\anonsection{СПИСОК ИСПОЛЬЗОВАННЫХ ИСТОЧНИКОВ}
\printbibliography[heading=none] % печать библиографии   title={СПИСОК ИСПОЛЬЗОВАННЫХ ИСТОЧНИКОВ}, heading=bibintoc


\newpage
% Приложения
\appsection{Приложение А} \hypertarget{app-a}{\label{app-a}}

\centering{Программный код файла FisheyeDewarper.hpp}
\begin{english}
%\lstinputlisting[language=C++, numbers=left]{code/FisheyeDewarper.hpp}
\end{english}

\clearpage


% \appsection{Приложение Б} \hypertarget{app-b}{\label{app-b}}

% \centering{Программный код файла FisheyeDewarper.cpp}
% \begin{english}
% \lstinputlisting[language=C++, numbers=left]{code/FisheyeDewarper.cpp}
% \end{english}

% \clearpage

% \appsection{Приложение В} \hypertarget{app-c}{\label{app-c}}

% \centering{Программный код файла unity\_plugin.cpp}
% \begin{english}
% \lstinputlisting[language=C++, numbers=left]{code/unity_plugin.cpp}
% \end{english}

% \clearpage

% \appsection{Приложение Г} \hypertarget{app-d}{\label{app-d}}

% \centering{Программный код файла Connector.cs}
% \begin{english}
% \lstinputlisting[language=C, numbers=left]{code/Connector.cs}
% \end{english}

% \clearpage

% \appsection{Приложение Д} \hypertarget{app:matlab}{\label{app:matlab}}

% \centering{Программный код файла DisparityPlayground.m}
% \begin{english}
% \lstinputlisting[language=Matlab, numbers=left]{code/DisparityPlayground.m}
% \end{english}

% \clearpage
 % Код скрипта

\end{document} 
%%% Конец документа
